\documentclass[a4paper,12pt]{article}

\usepackage[utf8]{inputenc}
\usepackage[portuguese]{babel}
\usepackage{fancyhdr}
\usepackage{lastpage}
\usepackage{amsmath}
\usepackage{amsthm}
\usepackage{amssymb}
\usepackage[portuguese,algoruled,lined,linesnumbered]{algorithm2e}
\usepackage{geometry}
\usepackage{hyperref}
\geometry{a4paper, margin=1in}


% ---- PREENCHA ----
\newcommand{\Lista}{5}
\newcommand{\Nome}{Leonardo Heidi Almeida Murakami}
\newcommand{\NUSP}{11260186}
% -----------------------------

\pagestyle{fancy}
\fancyhf{}
\fancyhead[L]{\Nome}
\fancyhead[R]{NUSP \NUSP}
\fancyfoot[L]{Lista \Lista}
\fancyfoot[C]{MAC0338}
\fancyfoot[R]{Página \thepage\ de \pageref{LastPage}}

\begin{document}

\section*{Exercício 7}

\subsection*{Estratégia do Algoritmo}
Usaremos PD. A ideia é, para cada posição $i$ do vetor, calcular o comprimento da maior subsequência crescente que termina especificamente em $v[i]$. Depois, basta pegar o máximo entre todos esses valores.

Definimos $L[i]$ como o comprimento da LIS (\textit{Longest Increasing Subsequence}) que \textbf{termina} no elemento $v[i]$. Note que isso não é a LIS do subvetor $v[1..i]$ inteiro — é especificamente a maior subsequência crescente que inclui $v[i]$ como último elemento.

A estratégia funciona porque qualquer LIS do vetor inteiro necessariamente termina em alguma posição $i$. Logo, se calcularmos corretamente $L[i]$ para todo $i$, o máximo de $L$ será a resposta.

Para calcular $L[i]$, aproveitamos os valores $L[j]$ já calculados para $j < i$. Se $v[j] < v[i]$, podemos "grudar" $v[i]$ no final da LIS que termina em $v[j]$, criando uma subsequência de tamanho $1 + L[j]$. Escolhemos o $j$ que maximiza esse valor.

\subsection*{Recorrência}
Formalizando a ideia acima:
\begin{equation}
L[i] = 1 + \max(\{L[j] \mid 1 \le j < i \text{ e } v[j] < v[i]\} \cup \{0\}),
\end{equation}
onde o $\max$ do conjunto vazio é 0.

A resposta final é o maior valor em $L$, já que a LIS pode terminar em qualquer posição:
\begin{equation}
\text{Resultado} = \max_{1 \le i \le n} L[i].
\end{equation}
\subsection*{Prova de Corretude}
Vamos provar por indução em $i$ que $L[i]$ é calculado corretamente.

\textbf{Base:} Para $i=1$, temos $L[1] = 1$. De fato, a única subsequência que termina em $v[1]$ é o próprio $v[1]$, que tem comprimento 1.

\textbf{Hipótese de Indução (HI):} Assuma que para todo $k < i$, o valor $L[k]$ está correto.

\textbf{Passo Indutivo:} Considere $i > 1$. A LIS terminando em $v[i]$ pode ser:
\begin{itemize}
    \item Uma LIS terminando em algum $v[j]$ (com $j < i$ e $v[j] < v[i]$), estendida por $v[i]$. Pela HI, $L[j]$ está correto, então o tamanho é $1 + L[j]$. Como queremos o máximo, escolhemos o maior $L[j]$ possível.
    \item Apenas $v[i]$ sozinho, se não existir $v[j] < v[i]$ antes dele. Neste caso o tamanho é 1, que a recorrência retorna como $1+0=1$.
\end{itemize}

A recorrência captura ambos os casos, logo $L[i]$ é correto. Como todo $L[i]$ está correto, o máximo de $L$ é a resposta certa.
\subsection*{Algoritmo}
\begin{algorithm}[H]
\SetAlgoLined
\caption{\texttt{SubsequenciaCrescenteMaisLonga}$(v, n)$}
\KwData{Vetor $v[1..n]$ de inteiros}
\KwResult{Comprimento da subsequência crescente mais longa de $v$}
\BlankLine
Seja $L[1..n]$ um novo vetor\;
\For{$i \leftarrow 1$ \KwTo $n$}{
    $L[i] \leftarrow 1$\;
}
\For{$i \leftarrow 2$ \KwTo $n$}{
    \For{$j \leftarrow 1$ \KwTo $i-1$}{
        \If{$v[j] < v[i]$}{
            $L[i] \leftarrow \max(L[i], 1 + L[j])$\;
        }
    }
}
comprimentoMax $\leftarrow 0$\;
\For{$i \leftarrow 1$ \KwTo $n$}{
    comprimentoMax $\leftarrow \max(comprimentoMax, L[i])$\;
}
\Return{comprimentoMax}\;
\end{algorithm}

\subsection*{Análise de Complexidade}
Seja $T(n)$ o tempo de execução do algoritmo. Analisamos cada parte:

\begin{enumerate}
    \item \textbf{Inicialização} (linhas 2-4): Inicializa $n$ posições do vetor $L$, levando tempo $O(n)$.
    
    \item \textbf{Cálculo de $L$} (linhas 5-11): Dois laços aninhados. O laço externo vai de $i=2$ até $n$. Para cada $i$, o laço interno vai de $j=1$ até $i-1$. O total de iterações é:
    \begin{align*}
    \sum_{i=2}^{n} (i-1) &= \sum_{k=1}^{n-1} k \\
    &= \frac{(n-1)n}{2} \\
    &= \frac{n^2 - n}{2}.
    \end{align*}
    Cada iteração faz operações $O(1)$, logo esta parte custa $O(n^2)$.
    
    \item \textbf{Busca do máximo} (linhas 12-15): Percorre $n$ elementos em $O(n)$.
\end{enumerate}

O custo total é dominado pelo cálculo de $L$:
\begin{equation}
\boxed{T(n) = O(n^2)}.
\end{equation}

O espaço usado é $O(n)$ para o vetor $L$.

\section*{Exercício 14}

\subsection*{Estratégia do Algoritmo}
Usaremos PD. A ideia é calcular o número de maneiras de formar cada valor $j$ (de 0 até $V$) usando as moedas disponíveis, adicionando uma de cada vez.

Seja $C[i][j]$ o número de composicoes para formar o valor $j$ usando apenas as $i$ primeiras denominações ($v[1], \dots, v[i]$), respeitando suas quantidades ($q[1], \dots, q[i]$).

\subsection*{Recorrência}
Para calcular $C[i][j]$, consideramos usar $k$ unidades da moeda $v[i]$, onde $k$ varia de $0$ até $q[i]$. Se usarmos $k$ unidades de $v[i]$, elas somam $k \cdot v[i]$. O resto, $j - k \cdot v[i]$, precisa ser formado com as primeiras $i-1$ denominações, o que pode ser feito de $C[i-1][j - k \cdot v[i]]$ maneiras.

Como não podemos usar simultaneamente quantidades diferentes de $v[i]$, somamos as possibilidades:
\begin{equation}
C[i][j] = \sum_{k=0}^{q[i]} C[i-1][j - k \cdot v[i]],
\end{equation}
onde apenas consideramos termos com $j - k \cdot v[i] \ge 0$.

\textbf{Casos Base:}
\begin{itemize}
    \item $C[0][0] = 1$: Uma maneira de formar 0 sem moedas (composição vazia).
    \item $C[0][j] = 0$ para $j > 0$: Impossível formar valor positivo sem moedas.
\end{itemize}

A resposta e $C[n][V]$.
\subsection*{Prova de Corretude}
Vamos provar por indução em $i$ que $C[i][j]$ está correto para todo $j$.

\textbf{Base:} Para $i=0$, temos $C[0][0]=1$ e $C[0][j]=0$ para $j>0$, que está correto.

\textbf{Hipótese de Indução (HI):} Assuma que $C[i-1][j]$ está correto para todo $j$.

\textbf{Passo Indutivo:} Para calcular $C[i][j]$, somamos as contagens para cada possível quantidade $k$ da moeda $v[i]$:
\begin{itemize}
    \item Usar 0 de $v[i]$: $C[i-1][j]$ maneiras
    \item Usar 1 de $v[i]$: $C[i-1][j-v[i]]$ maneiras
    \item ...
    \item Usar $k$ de $v[i]$: $C[i-1][j-k \cdot v[i]]$ maneiras
\end{itemize}

Pela HI, cada termo está correto. Como somamos todas as possibilidades (de $k=0$ até $k=q[i]$) e elas são disjuntas, temos $C[i][j]$ correto. Logo $C[n][V]$ é a resposta certa.

\subsection*{Algoritmo}
\begin{algorithm}[H]
\SetAlgoLined
\caption{\texttt{NumeroComposicoes}$(v, q, n, V)$}
\KwData{Vetores $v[1..n]$, $q[1..n]$, inteiros $n > 0$, $V > 0$}
\KwResult{O número de composições de $V$}
\BlankLine
Seja $C[0..n][0..V]$ uma nova matriz 2D, inicializada com zeros\;
$C[0][0] \leftarrow 1$\;
\For{$i \leftarrow 1$ \KwTo $n$}{
    \For{$j \leftarrow 0$ \KwTo $V$}{
        \For{$k \leftarrow 0$ \KwTo $q[i]$}{
            \If{$j - k \cdot v[i] \ge 0$}{
                $C[i][j] \leftarrow C[i][j] + C[i-1][j - k \cdot v[i]]$\;
            }
        }
    }
}
\Return{$C[n][V]$}\;
\end{algorithm}

\subsection*{Análise de Complexidade}
Temos três laços aninhados:
\begin{enumerate}
    \item Laço sobre as $n$ denominações (linha 3).
    \item Laço sobre os valores de $0$ até $V$ (linha 4): $V+1$ iterações.
    \item Laço sobre a quantidade $k$ de $0$ a $q[i]$ (linha 5): no pior caso $\max(q)+1$ vezes.
\end{enumerate}

O total de operações é:
\begin{align*}
T(n, V) &= \sum_{i=1}^{n} \sum_{j=0}^{V} \sum_{k=0}^{q[i]} O(1) \\
&= n \cdot V \cdot \max(q).
\end{align*}

Logo:
\begin{equation}
\boxed{O(n \cdot V \cdot \max(q))}.
\end{equation}

O espaço usado é $O(n \cdot V)$ para a matriz $C$.



\end{document}